
\Data\ID{1003102401}
\Updated{2023-04-03 01:04:40}
\Level{A}
\SubArea{Calculations with logarithms}
\LANGen{
\Question{
Among offered answers, choose a logarithmic form of the following equality.
\[ 8^2  = 64 \]}
{\( \log_864=2 \)}{\( \log_264=8 \)}{\( \log_82=64 \)}{\( \log_28=64 \)}{}{}}
\LANGpl{
\Question{
Która z poniższych opcji przedstawia podaną równość w formie logarytmicznej? 
\[ 8^2  = 64 \]}
{\( \log_864=2 \)}{\( \log_264=8 \)}{\( \log_82=64 \)}{\( \log_28=64 \)}{}{}}
\LANGcs{
\Question{
Zapište rovnost pomocí logaritmu.
\[ 8^2  = 64 \]}
{\( \log_864=2 \)}{\( \log_264=8 \)}{\( \log_82=64 \)}{\( \log_28=64 \)}{}{}}
\LANGsk{
\Question{
Zapíšte rovnosť v tvare logaritmu.
\[ 8^2  = 64 \]}
{\( \log_864=2 \)}{\( \log_264=8 \)}{\( \log_82=64 \)}{\( \log_28=64 \)}{}{}}
\LANGes{
\Question{
Expresa en forma logarítmica la siguiente igualdad. 
\[ 8^2  = 64 \]}
{\( \log_864=2 \)}{\( \log_264=8 \)}{\( \log_82=64 \)}{\( \log_28=64 \)}{}{}}
\End

\Data\ID{1003102402}
\Updated{2023-04-03 01:04:36}
\Level{A}
\SubArea{Calculations with logarithms}
\LANGen{
\Question{
Find the value of \( y \) if 
\[ \log_42=y. \]}
{\( y=\frac12 \)}{\( y=-1 \)}{\( y=2 \)}{\( y=1 \)}{}{}}
\LANGpl{
\Question{
Wyznacz wartość \( y \) jeśli
\[ \log_42=y. \]}
{\( y=\frac12 \)}{\( y=-1 \)}{\( y=2 \)}{\( y=1 \)}{}{}}
\LANGcs{
\Question{
Určete hodnotu čísla \( y \), aby platilo:
\[ \log_42=y. \]}
{\( y=\frac12 \)}{\( y=-1 \)}{\( y=2 \)}{\( y=1 \)}{}{}}
\LANGsk{
\Question{
Určte hodnotu čísla \( y \), aby platilo:
\[ \log_42=y. \]}
{\( y=\frac12 \)}{\( y=-1 \)}{\( y=2 \)}{\( y=1 \)}{}{}}
\LANGes{
\Question{
Halla el valor de \( y \) sabiendo que :
\[ \log_42=y. \]}
{\( y=\frac12 \)}{\( y=-1 \)}{\( y=2 \)}{\( y=1 \)}{}{}}
\End

\Data\ID{1003102403}
\Updated{2023-04-03 01:04:56}
\Level{A}
\SubArea{Calculations with logarithms}
\LANGen{
\Question{
Find the value of \( x \) if
\[ \log_x8=-3. \]}
{\( x=\frac12 \)}{\( x=2 \)}{\( x=\sqrt2 \)}{\( x=-\frac12 \)}{}{}}
\LANGpl{
\Question{
Wyznacz wartość \( x \) jeśli
\[ \log_x8=-3. \]}
{\( x=\frac12 \)}{\( x=2 \)}{\( x=\sqrt2 \)}{\( x=-\frac12 \)}{}{}}
\LANGcs{
\Question{
Určete hodnotu čísla \( x \), aby platilo:
\[ \log_x8=-3. \]}
{\( x=\frac12 \)}{\( x=2 \)}{\( x=\sqrt2 \)}{\( x=-\frac12 \)}{}{}}
\LANGsk{
\Question{
Určte hodnotu čísla \( x \), aby platilo:
\[ \log_x8=-3. \]}
{\( x=\frac12 \)}{\( x=2 \)}{\( x=\sqrt2 \)}{\( x=-\frac12 \)}{}{}}
\LANGes{
\Question{
Halla el valor de \( x \) sabiendo que :
\[ \log_x8=-3. \]}
{\( x=\frac12 \)}{\( x=2 \)}{\( x=\sqrt2 \)}{\( x=-\frac12 \)}{}{}}
\End

\Data\ID{1003102404}
\Updated{2023-04-03 01:04:12}
\Level{A}
\SubArea{Calculations with logarithms}
\LANGen{
\Question{
Which of the following statements is  not  true?}
{\( \log_{\sqrt2}2=\frac12 \)}{\( \log_{\frac12}2=-1 \)}{\( \log_21= 0 \)}{\( \log_2\sqrt2=\frac12 \)}{}{}}
\LANGpl{
\Question{
Które z poniższych stwierdzeń  nie jest prawdziwe?}
{\( \log_{\sqrt2}2=\frac12 \)}{\( \log_{\frac12}2=-1 \)}{\( \log_21= 0 \)}{\( \log_2\sqrt2=\frac12 \)}{}{}}
\LANGcs{
\Question{
Který z výroků neplatí?}
{\( \log_{\sqrt2}2=\frac12 \)}{\( \log_{\frac12}2=-1 \)}{\( \log_21= 0 \)}{\( \log_2\sqrt2=\frac12 \)}{}{}}
\LANGsk{
\Question{
Ktorý z nasledujúcich výrokov  nie  je pravdivý?}
{\( \log_{\sqrt2}2=\frac12 \)}{\( \log_{\frac12}2=-1 \)}{\( \log_21= 0 \)}{\( \log_2\sqrt2=\frac12 \)}{}{}}
\LANGes{
\Question{
¿Cuál de los siguientes enunciados  no  es verdadero?}
{\( \log_{\sqrt2}2=\frac12 \)}{\( \log_{\frac12}2=-1 \)}{\( \log_21= 0 \)}{\( \log_2\sqrt2=\frac12 \)}{}{}}
\End

\Data\ID{1003102405}
\Updated{2023-04-03 01:04:24}
\Level{A}
\SubArea{Calculations with logarithms}
\LANGen{
\Question{
If \( \log_x\frac49=-\frac12 \) then \( x \) is equal to:}
{\( \frac{81}{16} \)}{\( \frac{16}{81} \)}{\( \frac23 \)}{\( -\frac23 \)}{}{}}
\LANGpl{
\Question{
Jeśli \( \log_x\frac49=-\frac12 \) wtedy \( x \) równe jest:}
{\( \frac{81}{16} \)}{\( \frac{16}{81} \)}{\( \frac23 \)}{\( -\frac23 \)}{}{}}
\LANGcs{
\Question{
Určete hodnotu čísla \( x \), aby platilo \( \log_x\frac49=-\frac12 \).}
{\( \frac{81}{16} \)}{\( \frac{16}{81} \)}{\( \frac23 \)}{\( -\frac23 \)}{}{}}
\LANGsk{
\Question{
Určte hodnotu čísla \( x \), aby platilo \( \log_x\frac49=-\frac12 \).}
{\( \frac{81}{16} \)}{\( \frac{16}{81} \)}{\( \frac23 \)}{\( -\frac23 \)}{}{}}
\LANGes{
\Question{
Si \( \log_x\frac49=-\frac12 \) entonces \( x \) es igual a:}
{\( \frac{81}{16} \)}{\( \frac{16}{81} \)}{\( \frac23 \)}{\( -\frac23 \)}{}{}}
\End

\Data\ID{1003102406}
\Updated{2023-04-03 01:04:46}
\Level{A}
\SubArea{Calculations with logarithms}
\LANGen{
\Question{
Find the value of \( z \), if
\( \log_3z=-3 \).}
{\( z=\frac1{27} \)}{\( z=27 \)}{\( z=-9 \)}{\( z=\frac19 \)}{}{}}
\LANGpl{
\Question{
Określ wartość \( z \) jeśli
\( \log_3z=-3 \).}
{\( z=\frac1{27} \)}{\( z=27 \)}{\( z=-9 \)}{\( z=\frac19 \)}{}{}}
\LANGcs{
\Question{
Určete hodnotu \( z \), tak, aby platilo
\( \log_3z=-3 \).}
{\( z=\frac1{27} \)}{\( z=27 \)}{\( z=-9 \)}{\( z=\frac19 \)}{}{}}
\LANGsk{
\Question{
Určte hodnotu čísla \( z \), aby platilo
\( \log_3z=-3 \).}
{\( z=\frac1{27} \)}{\( z=27 \)}{\( z=-9 \)}{\( z=\frac19 \)}{}{}}
\LANGes{
\Question{
Halla el valor de \( z \) si
\( \log_3z=-3 \).}
{\( z=\frac1{27} \)}{\( z=27 \)}{\( z=-9 \)}{\( z=\frac19 \)}{}{}}
\End

\Data\ID{1003102408}
\Updated{2023-04-03 02:04:10}
\Level{A}
\SubArea{Calculations with logarithms}
\LANGen{
\Question{
The value of the expression \( \log_{16}\left(\log_216 \right) \) is:}
{\( \frac12 \)}{\( \frac14 \)}{\( \frac18 \)}{\( 2 \)}{}{}}
\LANGpl{
\Question{
Wartość wyrażenia  \( \log_{16}\left(\log_216 \right) \) wynosi:}
{\( \frac12 \)}{\( \frac14 \)}{\( \frac18 \)}{\( 2 \)}{}{}}
\LANGcs{
\Question{
Hodnota výrazu \( \log_{16}\left(\log_216 \right) \) se rovná číslu:}
{\( \frac12 \)}{\( \frac14 \)}{\( \frac18 \)}{\( 2 \)}{}{}}
\LANGsk{
\Question{
Hodnota výrazu \( \log_{16}\left(\log_216 \right) \) sa rovná číslu:}
{\( \frac12 \)}{\( \frac14 \)}{\( \frac18 \)}{\( 2 \)}{}{}}
\LANGes{
\Question{
El valor de la expresión \( \log_{16}\left(\log_216 \right) \) es:}
{\( \frac12 \)}{\( \frac14 \)}{\( \frac18 \)}{\( 2 \)}{}{}}
\End

\Data\ID{2000000606}
\Updated{2023-04-03 01:04:37}
\Level{A}
\SubArea{Calculations with logarithms}
\LANGen{
\Question{
The value of the expression \(3\log_{4}2 - \frac{1}{2}\log_{4}16\) is:}
{\(0.5\)}{\(2\)}{\(\log_{4}\frac{9}{256}\)}{\( \frac{3}{2} \log_{4}\frac{1}{8}\)}{}{}}
\LANGpl{
\Question{
Wartością wyrażenia  \(3\log_{4}2 - \frac{1}{2}\log_{4}16\) jest:}
{\(0{,}5\)}{\(2\)}{\(\log_{4}\frac{9}{256}\)}{\( \frac{3}{2} \log_{4}\frac{1}{8}\)}{}{}}
\LANGcs{
\Question{
Hodnota výrazu \(3\log_{4}2 - \frac{1}{2}\log_{4}16\) je:}
{\(0{,}5\)}{\(2\)}{\(\log_{4}\frac{9}{256}\)}{\( \frac{3}{2} \log_{4}\frac{1}{8}\)}{}{}}
\LANGsk{
\Question{
Hodnota výrazu \(3\log_{4}2 - \frac{1}{2}\log_{4}16\) je:}
{\(0{,}5\)}{\(2\)}{\(\log_{4}\frac{9}{256}\)}{\( \frac{3}{2} \log_{4}\frac{1}{8}\)}{}{}}
\LANGes{
\Question{
El valor de la expresión \(3\log_{4}2 - \frac{1}{2}\log_{4}16\) es:}
{\(0.5\)}{\(2\)}{\(\log_{4}\frac{9}{256}\)}{\( \frac{3}{2} \log_{4}\frac{1}{8}\)}{}{}}
\End

\Data\ID{2000014110}
\Updated{2023-04-03 01:04:33}
\Level{A}
\SubArea{Calculations with logarithms}
\LANGen{
\Question{
Identify which of the following equalities is not equivalent to the equality \(4^x=9\).}
{\( x=2\log_2 9\)}{\( x=\log_2 3\)}{\( x=\log_4 9\)}{\( x=2\log_4 3\)}{}{}}
\LANGpl{
\Question{
Wskaż równość, która nie jest równa równaniu \(4^x=9\).}
{\( x=2\log_2 9\)}{\( x=\log_2 3\)}{\( x=\log_4 9\)}{\( x=2\log_4 3\)}{}{}}
\LANGcs{
\Question{
Rozhodněte, která z následujících rovnic není ekvivalentní s rovnicí \(4^x=9\).}
{\( x=2\log_2 9\)}{\( x=\log_2 3\)}{\( x=\log_4 9\)}{\( x=2\log_4 3\)}{}{}}
\LANGsk{
\Question{
Rozhodnite, ktorá z nasledujúcich rovníc nie je ekvivalentná s rovnicou  \(4^x=9\).}
{\( x=2\log_2 9\)}{\( x=\log_2 3\)}{\( x=\log_4 9\)}{\( x=2\log_4 3\)}{}{}}
\LANGes{
\Question{
Identifica cuál de las siguientes igualdades no es equivalente a la igualdad \(4^x=9\).}
{\( x=2\log_2 9\)}{\( x=\log_2 3\)}{\( x=\log_4 9\)}{\( x=2\log_4 3\)}{}{}}
\End

\Data\ID{2010011001}
\Updated{2023-04-03 01:04:04}
\Level{A}
\SubArea{Calculations with logarithms}
\LANGen{
\Question{
Among offered answers, choose a logarithmic form of the following equality.
\[ \sqrt{16}  = 4 \]}
{\( \log_{16}4=\frac{1}{2}\)}{\( \log_{\frac12}16=4\)}{\( \log_4 \frac12=16\)}{\( \log_{2}4=16\)}{}{}}
\LANGpl{
\Question{
Napisz równość w postaci logarytmu.
\[ \sqrt{16}  = 4 \]}
{\( \log_{16}4=\frac{1}{2}\)}{\( \log_{\frac12}16=4\)}{\( \log_4 \frac12=16\)}{\( \log_{2}4=16\)}{}{}}
\LANGcs{
\Question{
Zapište rovnost ve tvaru logaritmu.
\[ \sqrt{16}  = 4 \]}
{\( \log_{16}4=\frac{1}{2}\)}{\( \log_{\frac12}16=4\)}{\( \log_4 \frac12=16\)}{\( \log_{2}4=16\)}{}{}}
\LANGsk{
\Question{
Zapíšte rovnosť v tvare logaritmu.
\[ \sqrt{16}  = 4 \]}
{\( \log_{16}4=\frac{1}{2}\)}{\( \log_{\frac12}16=4\)}{\( \log_4 \frac12=16\)}{\( \log_{2}4=16\)}{}{}}
\LANGes{
\Question{
Entre las siguientes respuestas, elige la forma logarítmica que representa la siguiente igualdad.
\[ \sqrt{16}  = 4 \]}
{\( \log_{16}4=\frac{1}{2}\)}{\( \log_{\frac12}16=4\)}{\( \log_4 \frac12=16\)}{\( \log_{2}4=16\)}{}{}}
\End

\Data\ID{2010011002}
\Updated{2023-04-03 01:04:26}
\Level{A}
\SubArea{Calculations with logarithms}
\LANGen{
\Question{
Which of the following statements is  not  true?}
{\( \log_{\frac12}6=-3 \)}{\( \log_{\frac12}8=-3\)}{\( \log_2 \sqrt{2}=\frac12\)}{\( \log_{\frac12}\frac14=2\)}{}{}}
\LANGpl{
\Question{
Które z poniższych stwierdzeń  nie  jest prawdziwe?}
{\( \log_{\frac12}6=-3 \)}{\( \log_{\frac12}8=-3\)}{\( \log_2 \sqrt{2}=\frac12\)}{\( \log_{\frac12}\frac14=2\)}{}{}}
\LANGcs{
\Question{
Které z následujících tvrzení není pravdivé?}
{\( \log_{\frac12}6=-3 \)}{\( \log_{\frac12}8=-3\)}{\( \log_2 \sqrt{2}=\frac12\)}{\( \log_{\frac12}\frac14=2\)}{}{}}
\LANGsk{
\Question{
Ktorý z nasledujúcich výrokov   nie   je pravdivý?}
{\( \log_{\frac12}6=-3 \)}{\( \log_{\frac12}8=-3\)}{\( \log_2 \sqrt{2}=\frac12\)}{\( \log_{\frac12}\frac14=2\)}{}{}}
\LANGes{
\Question{
¿Cuál de las siguientes afirmaciones  no  es verdadera?}
{\( \log_{\frac12}6=-3 \)}{\( \log_{\frac12}8=-3\)}{\( \log_2 \sqrt{2}=\frac12\)}{\( \log_{\frac12}\frac14=2\)}{}{}}
\End

\Data\ID{2010011003}
\Updated{2023-04-03 01:04:43}
\Level{A}
\SubArea{Calculations with logarithms}
\LANGen{
\Question{
Find the value of \( x \), if
\( \log_{\frac13}x=-4 \).}
{\( x=81 \)}{\( x=\frac1{81} \)}{\( x=-81 \)}{\( x=\frac1{12} \)}{}{}}
\LANGpl{
\Question{
Znajdź wartość \( x \), jeśli
\( \log_{\frac13}x=-4 \).}
{\( x=81 \)}{\( x=\frac1{81} \)}{\( x=-81 \)}{\( x=\frac1{12} \)}{}{}}
\LANGcs{
\Question{
Určete hodnotu \( x \) tak, aby
\( \log_{\frac13}x=-4 \).}
{\( x=81 \)}{\( x=\frac1{81} \)}{\( x=-81 \)}{\( x=\frac1{12} \)}{}{}}
\LANGsk{
\Question{
Určte hodnotu čísla \( x \), aby platilo \( \log_{\frac13}x=-4 \).}
{\( x=81 \)}{\( x=\frac1{81} \)}{\( x=-81 \)}{\( x=\frac1{12} \)}{}{}}
\LANGes{
\Question{
Halla el valor de \( x \), si
\( \log_{\frac13}x=-4 \).}
{\( x=81 \)}{\( x=\frac1{81} \)}{\( x=-81 \)}{\( x=\frac1{12} \)}{}{}}
\End

\Data\ID{2010016007}
\Updated{2023-04-03 01:04:35}
\Level{A}
\SubArea{Calculations with logarithms}
\LANGen{
\Question{
If \(\log_3 a=b\) then the value of \(\log_9 a\) is equal to:}
{\(\frac{b}2\)}{\(2b\)}{\( \frac2{b}\)}{\(9b\)}{}{}}
\LANGpl{
\Question{
Jeżeli \(\log_3 a=b\), to wartość \(\log_9 a\) jest równa:}
{\(\frac{b}2\)}{\(2b\)}{\( \frac2{b}\)}{\(9b\)}{}{}}
\LANGcs{
\Question{
Když \(\log_3 a=b\), pak je hodnota \(\log_9 a\) rovna:}
{\(\frac{b}2\)}{\(2b\)}{\( \frac2{b}\)}{\(9b\)}{}{}}
\LANGsk{
\Question{
Ak \(\log_3 a=b\), tak hodnota \(\log_9 a\) sa rovná:}
{\(\frac{b}2\)}{\(2b\)}{\( \frac2{b}\)}{\(9b\)}{}{}}
\LANGes{
\Question{
Si \(\log_3 a=b\) entonces el valor de \(\log_9 a\) es igual a:}
{\(\frac{b}2\)}{\(2b\)}{\( \frac2{b}\)}{\(9b\)}{}{}}
\End

\Data\ID{2010016008}
\Updated{2023-04-03 01:04:06}
\Level{A}
\SubArea{Calculations with logarithms}
\LANGen{
\Question{
If \(\log_2 a=b\) then the value of \(\log_8 a\) is equal to:}
{\(\frac{b}3\)}{\(\frac{b}2\)}{\(3b\)}{\(4b\)}{}{}}
\LANGpl{
\Question{
Jeżeli \(\log_2 a=b\), to wartość  \(\log_8 a\) jest równa:}
{\(\frac{b}3\)}{\(\frac{b}2\)}{\(3b\)}{\(4b\)}{}{}}
\LANGcs{
\Question{
Když \(\log_2 a=b\), pak je hodnota \(\log_8 a\) rovna:}
{\(\frac{b}3\)}{\(\frac{b}2\)}{\(3b\)}{\(4b\)}{}{}}
\LANGsk{
\Question{
Ak \(\log_2 a=b\), tak hodnota \(\log_8 a\) sa rovná:}
{\(\frac{b}3\)}{\(\frac{b}2\)}{\(3b\)}{\(4b\)}{}{}}
\LANGes{
\Question{
Si \(\log_2 a=b\) entonces el valor de \(\log_8 a\) es igual a:}
{\(\frac{b}3\)}{\(\frac{b}2\)}{\(3b\)}{\(4b\)}{}{}}
\End

\Data\ID{9000022802}
\Updated{2023-04-03 02:04:13}
\Level{A}
\SubArea{Calculations with logarithms}
\LANGen{
\Question{
Find all \(x\in \mathbb{R}\) 
for which the following expression is undefined. 
\[ 
 \log \left (2x^{2} + 4x - 6\right )
\]}
{\(\left [ -3;1\right ] \)}{\(\left (-\infty ;-3\right )\cup \left (1;\infty \right )\)}{\(\left (-3;1\right )\)}{\(\left (-\infty ;-3\right ] \cup \left [ 1;\infty \right )\)}{}{}}
\LANGpl{
\Question{
Wyznacz wszystkie wartości \(x\in \mathbb{R}\), dla których podane wyrażenie nie jest określone.
\[ 
 \log \left (2x^{2} + 4x - 6\right )
\]}
{\(\left [ -3;1\right ] \)}{\(\left (-\infty ;-3\right )\cup \left (1;\infty \right )\)}{\(\left (-3;1\right )\)}{\(\left (-\infty ;-3\right ] \cup \left [ 1;\infty \right )\)}{}{}}
\LANGcs{
\Question{
Množina všech \(x\in \mathbb{R}\), pro která
není definován výraz \(\log \left (2x^{2} + 4x - 6\right )\),
je:}
{\(\left \langle -3;1\right \rangle \)}{\(\left (-\infty ;-3\right )\cup \left (1;\infty \right )\)}{\(\left (-3;1\right )\)}{\(\left (-\infty ;-3\right \rangle \cup \left \langle 1;\infty \right )\)}{}{}}
\LANGsk{
\Question{
Množina všetkých \(x\in \mathbb{R}\), pre ktoré
nie je definovaný výraz \(\log \left (2x^{2} + 4x - 6\right )\),
je:}
{\(\left \langle -3;1\right \rangle \)}{\(\left (-\infty ;-3\right )\cup \left (1;\infty \right )\)}{\(\left (-3;1\right )\)}{\(\left (-\infty ;-3\right \rangle \cup \left \langle 1;\infty \right )\)}{}{}}
\LANGes{
\Question{
Halla el conjunto de todos los \(x\in \mathbb{R}\) 
para los que la siguiente expresión no está definida. 
\[ 
 \log \left (2x^{2} + 4x - 6\right )
\]}
{\(\left [ -3;1\right ] \)}{\(\left (-\infty ;-3\right )\cup \left (1;\infty \right )\)}{\(\left (-3;1\right )\)}{\(\left (-\infty ;-3\right ] \cup \left [ 1;\infty \right )\)}{}{}}
\End

\Data\ID{9000034902}
\Updated{2023-04-03 01:04:01}
\Level{A}
\SubArea{Calculations with logarithms}
\LANGen{
\Question{
Find the domain of the following expression. 
\[ 
 \log _{2}\left [\left (\frac{2}
{3} - x\right )\left (x + \frac{1} 
{4}\right )\right ]
\]}
{\(\left (-\frac{1}
{4}; \frac{2} 
{3}\right )\)}{\(\left (-\infty ;-\frac{1}
{4}\right ] \cup \left [ \frac{2}
{3};\infty \right )\)}{\(\left (-\infty ;-\frac{1}
{4}\right )\cup \left (\frac{2}
{3};\infty \right )\)}{\(\left [ \frac{1}
{4}; \frac{2} 
{3}\right ] \)}{}{}}
\LANGpl{
\Question{
Wyznacz dziedzinę następującego wyrażenia.
\[ 
 \log _{2}\left [\left (\frac{2}
{3} - x\right )\left (x + \frac{1} 
{4}\right )\right]
\]}
{\(\left (-\frac{1}
{4}; \frac{2} 
{3}\right )\)}{\(\left (-\infty ;-\frac{1}
{4}\right ] \cup \left [ \frac{2}
{3};\infty \right )\)}{\(\left (-\infty ;-\frac{1}
{4}\right )\cup \left (\frac{2}
{3};\infty \right )\)}{\(\left [ \frac{1}
{4}; \frac{2} 
{3}\right ] \)}{}{}}
\LANGcs{
\Question{
Množina všech \(x\in \mathbb{R}\), pro
která je definován výraz \(\log _{2}\left [\left (\frac{2}
{3} - x\right )\left (x + \frac{1} 
{4}\right )\right]\!,\) 
je:}
{\(\left (-\frac{1}
{4}; \frac{2} 
{3}\right )\)}{\(\left (-\infty ;-\frac{1}
{4}\right \rangle \cup \left \langle \frac{2}
{3};\infty \right )\)}{\(\left (-\infty ;-\frac{1}
{4}\right )\cup \left (\frac{2}
{3};\infty \right )\)}{\(\left \langle \frac{1}
{4}; \frac{2} 
{3}\right \rangle \)}{}{}}
\LANGsk{
\Question{
Nájdite definičný obor daného výrazu \(\log _{2}\left [\left (\frac{2}
{3} - x\right )\left (x + \frac{1} 
{4}\right )\right ]\).}
{\(\left (-\frac{1}
{4}; \frac{2} 
{3}\right )\)}{\(\left (-\infty ;-\frac{1}
{4}\right \rangle \cup \left \langle \frac{2}
{3};\infty \right )\)}{\(\left (-\infty ;-\frac{1}
{4}\right )\cup \left (\frac{2}
{3};\infty \right )\)}{\(\left \langle \frac{1}
{4}; \frac{2} 
{3}\right \rangle \)}{}{}}
\LANGes{
\Question{
Halla el dominio de la siguiente expresión
\[ 
 \log _{2}\left [\left (\frac{2}
{3} - x\right )\left (x + \frac{1} 
{4}\right )\right ]
\]}
{\(\left (-\frac{1}
{4}; \frac{2} 
{3}\right )\)}{\(\left (-\infty ;-\frac{1}
{4}\right ] \cup \left [ \frac{2}
{3};\infty \right )\)}{\(\left (-\infty ;-\frac{1}
{4}\right )\cup \left (\frac{2}
{3};\infty \right )\)}{\(\left [ \frac{1}
{4}; \frac{2} 
{3}\right ] \)}{}{}}
\End

\Data\ID{9000034904}
\Updated{2023-04-03 01:04:18}
\Level{A}
\SubArea{Calculations with logarithms}
\LANGen{
\Question{
Find all \(x\in \mathbb{R}\) 
for which the following expression is undefined. 
\[ 
 \log _{\frac{1} 
{4} }\left [\left (x + \frac{1} 
{2}\right )\left (5 - 2x\right )\right ]
\]}
{\(\left (-\infty ;-\frac{1}
{2}\right ] \cup \left [ \frac{5}
{2};\infty \right )\)}{\(\left [ -\frac{1}
{2}; \frac{5} 
{2}\right ] \)}{\(\left (-\frac{1}
{2}; \frac{5} 
{2}\right )\)}{\(\left (-\infty ;-\frac{1}
{2}\right )\cup \left (\frac{5}
{2};\infty \right )\)}{}{}}
\LANGpl{
\Question{
Wyznacz wszystkie wartości \(x\in \mathbb{R}\), dla których podane wyrażenie nie jest określone. 
\[ 
 \log _{\frac{1} 
{4} }\left [\left (x + \frac{1} 
{2}\right )\left (5 - 2x\right )\right ]
\]}
{\(\left (-\infty ;-\frac{1}
{2}\right ] \cup \left [ \frac{5}
{2};\infty \right )\)}{\(\left [ -\frac{1}
{2}; \frac{5} 
{2}\right ] \)}{\(\left (-\frac{1}
{2}; \frac{5} 
{2}\right )\)}{\(\left (-\infty ;-\frac{1}
{2}\right )\cup \left (\frac{5}
{2};\infty \right )\)}{}{}}
\LANGcs{
\Question{
Množina všech \(x\in \mathbb{R}\), pro která
není definován výraz \(\log _{\frac{1} 
{4} }\left [\left (x + \frac{1} 
{2}\right )\left (5 - 2x\right )\right]\!,\) 
je:}
{\(\left (-\infty ;-\frac{1}
{2}\right \rangle \cup \left \langle \frac{5}
{2};\infty \right )\)}{\(\left \langle -\frac{1}
{2}; \frac{5} 
{2}\right \rangle \)}{\(\left (-\frac{1}
{2}; \frac{5} 
{2}\right )\)}{\(\left (-\infty ;-\frac{1}
{2}\right )\cup \left (\frac{5}
{2};\infty \right )\)}{}{}}
\LANGsk{
\Question{
Množina všetkých \(x\in \mathbb{R}\), pre ktoré nie je definovaný výraz \(\log _{\frac{1} 
{4} }\left [\left (x + \frac{1} 
{2}\right )\left (5 - 2x\right )\right]\!,\) 
je:}
{\(\left (-\infty ;-\frac{1}
{2}\right \rangle \cup \left \langle \frac{5}
{2};\infty \right )\)}{\(\left \langle -\frac{1}
{2}; \frac{5} 
{2}\right \rangle \)}{\(\left (-\frac{1}
{2}; \frac{5} 
{2}\right )\)}{\(\left (-\infty ;-\frac{1}
{2}\right )\cup \left (\frac{5}
{2};\infty \right )\)}{}{}}
\LANGes{
\Question{
Halla todos los \(x\in \mathbb{R}\) 
para los cuales no está definida la siguiente expresión. 
\[ 
 \log _{\frac{1} 
{4} }\left [\left (x + \frac{1} 
{2}\right )\left (5 - 2x\right )\right ]
\]}
{\(\left (-\infty ;-\frac{1}
{2}\right ] \cup \left [ \frac{5}
{2};\infty \right )\)}{\(\left [ -\frac{1}
{2}; \frac{5} 
{2}\right ] \)}{\(\left (-\frac{1}
{2}; \frac{5} 
{2}\right )\)}{\(\left (-\infty ;-\frac{1}
{2}\right )\cup \left (\frac{5}
{2};\infty \right )\)}{}{}}
\End

\Data\ID{1003102409}
\Updated{2023-04-03 02:04:24}
\Level{B}
\SubArea{Calculations with logarithms}
\LANGen{
\Question{
Without a calculator evaluate the following expression and select the correct value.
\[ \log_64+\log_69 \]}
{\( 2 \)}{\( 36 \)}{\( 13 \)}{\( 6 \)}{}{}}
\LANGpl{
\Question{
Nie używając kalkulatora oszacuj podane wyrażenie i wybierz poprawną wartość. 
\[ \log_64+\log_69 \]}
{\( 2 \)}{\( 36 \)}{\( 13 \)}{\( 6 \)}{}{}}
\LANGcs{
\Question{
Bez použití kalkulačky vypočítejte hodnotu daného výrazu.
\[ \log_64+\log_69 \]}
{\( 2 \)}{\( 36 \)}{\( 13 \)}{\( 6 \)}{}{}}
\LANGsk{
\Question{
Bez použitia kalkulačky vypočítajte hodnotu daného výrazu.
\[ \log_64+\log_69 \]}
{\( 2 \)}{\( 36 \)}{\( 13 \)}{\( 6 \)}{}{}}
\LANGes{
\Question{
Sin calculadora, evalúa la siguiente expresión y elige el valor correcto. 
\[ \log_64+\log_69 \]}
{\( 2 \)}{\( 36 \)}{\( 13 \)}{\( 6 \)}{}{}}
\End

\Data\ID{1003102411}
\Updated{2023-04-03 01:04:54}
\Level{B}
\SubArea{Calculations with logarithms}
\LANGen{
\Question{
Without a calculator evaluate the following expression and select the correct value.
\[ \frac{\log\sqrt6}{\log6} \]}
{\( \frac12 \)}{\( 2 \)}{\( \frac{\sqrt6}6 \)}{\( \sqrt6-6 \)}{}{}}
\LANGpl{
\Question{
Nie używając kalkulatora oszacuj wyrażenie i wybierz poprawną wartość. 
\[ \frac{\log\sqrt6}{\log6} \]}
{\( \frac12 \)}{\( 2 \)}{\( \frac{\sqrt6}6 \)}{\( \sqrt6-6 \)}{}{}}
\LANGcs{
\Question{
Bez použití kalkulačky vypočítejte hodnotu daného výrazu.
\[ \frac{\log\sqrt6}{\log6} \]}
{\( \frac12 \)}{\( 2 \)}{\( \frac{\sqrt6}6 \)}{\( \sqrt6-6 \)}{}{}}
\LANGsk{
\Question{
Bez použitia kalkulačky vypočítajte hodnotu daného výrazu.
\[ \frac{\log\sqrt6}{\log6} \]}
{\( \frac12 \)}{\( 2 \)}{\( \frac{\sqrt6}6 \)}{\( \sqrt6-6 \)}{}{}}
\LANGes{
\Question{
Sin calculadora, evalúa la siguiente expresión y elige el valor correcto. 
\[ \frac{\log\sqrt6}{\log6} \]}
{\( \frac12 \)}{\( 2 \)}{\( \frac{\sqrt6}6 \)}{\( \sqrt6-6 \)}{}{}}
\End

\Data\ID{1003102412}
\Updated{2023-04-03 01:04:12}
\Level{B}
\SubArea{Calculations with logarithms}
\LANGen{
\Question{
If \( a \), \( b \), \( c\in(0;\infty) \) then the expression \( \log_5a-\frac23 \log_5 b+3\log_5c \) is equivalent to:}
{\( \log_5\frac{ac^3}{\sqrt[3]{b^2}} \)}{\( \log_5\frac{a\sqrt[3]{b^2}}{c^3} \)}{\( \log_5\frac{3ac}{\frac23 b} \)}{\( \log_5\frac{\frac23 ab}{3c} \)}{}{}}
\LANGpl{
\Question{
Jeśli  \( a \), \( b \), \( c\in(0;\infty) \), wtedy wyrażenie \( \log_5a-\frac23 \log_5 b+3\log_5c \) równe jest:}
{\( \log_5\frac{ac^3}{\sqrt[3]{b^2}} \)}{\( \log_5\frac{a\sqrt[3]{b^2}}{c^3} \)}{\( \log_5\frac{3ac}{\frac23 b} \)}{\( \log_5\frac{\frac23 ab}{3c} \)}{}{}}
\LANGcs{
\Question{
Upravte na jeden logaritmus výraz \( \log_5a-\frac23 \log_5 b+3\log_5c \), platí-li  \( a \), \( b \), \( c\in(0;\infty) \).}
{\( \log_5\frac{ac^3}{\sqrt[3]{b^2}} \)}{\( \log_5\frac{a\sqrt[3]{b^2}}{c^3} \)}{\( \log_5\frac{3ac}{\frac23 b} \)}{\( \log_5\frac{\frac23 ab}{3c} \)}{}{}}
\LANGsk{
\Question{
Upravte na jeden logaritmus výraz \( \log_5a-\frac23 \log_5 b+3\log_5c \), ak  \( a \), \( b \), \( c\in(0;\infty) \).}
{\( \log_5\frac{ac^3}{\sqrt[3]{b^2}} \)}{\( \log_5\frac{a\sqrt[3]{b^2}}{c^3} \)}{\( \log_5\frac{3ac}{\frac23 b} \)}{\( \log_5\frac{\frac23 ab}{3c} \)}{}{}}
\LANGes{
\Question{
Si \( a \), \( b \), \( c\in(0;\infty) \), la expresión \( \log_5a-\frac23 \log_5 b+3\log_5c \) es equivalente a:}
{\( \log_5\frac{ac^3}{\sqrt[3]{b^2}} \)}{\( \log_5\frac{a\sqrt[3]{b^2}}{c^3} \)}{\( \log_5\frac{3ac}{\frac23 b} \)}{\( \log_5\frac{\frac23 ab}{3c} \)}{}{}}
\End

\Data\ID{1003102413}
\Updated{2023-04-03 01:04:28}
\Level{B}
\SubArea{Calculations with logarithms}
\LANGen{
\Question{
Let \( x \), \( y \), \( z\in (0;\infty) \). Find the equivalent form of the following expression.
\[ \log\sqrt{\frac{xz^2}{y^{16}}} \]}
{\( \frac12\log x-8\log y+\log z \)}{\( \frac12\log x+8\log y-\log z \)}{\( 8\log x+\frac12\log y-\log z \)}{\( \log x-16\log y+2\log z \)}{}{}}
\LANGpl{
\Question{
Niech \( x \), \( y \), \( z\in (0;\infty) \). Wyznacz odpowiednik podanego wyrażenia.
\[ \log\sqrt{\frac{xz^2}{y^{16}}} \]}
{\( \frac12\log x-8\log y+\log z \)}{\( \frac12\log x+8\log y-\log z \)}{\( 8\log x+\frac12\log y-\log z \)}{\( \log x-16\log y+2\log z \)}{}{}}
\LANGcs{
\Question{
Nechť \( x \), \( y \), \( z\in (0;\infty) \). Daný výraz \[ \log\sqrt{\frac{xz^2}{y^{16}}} \] se rovná:}
{\( \frac12\log x-8\log y+\log z \)}{\( \frac12\log x+8\log y-\log z \)}{\( 8\log x+\frac12\log y-\log z \)}{\( \log x-16\log y+2\log z \)}{}{}}
\LANGsk{
\Question{
Nech \( x \), \( y \), \( z\in (0;\infty) \). Daný výraz \[ \log\sqrt{\frac{xz^2}{y^{16}}} \] sa rovná:}
{\( \frac12\log x-8\log y+\log z \)}{\( \frac12\log x+8\log y-\log z \)}{\( 8\log x+\frac12\log y-\log z \)}{\( \log x-16\log y+2\log z \)}{}{}}
\LANGes{
\Question{
Sea \( x \), \( y \), \( z\in (0;\infty) \). Halla la fórmula equivalente de la siguiente expresión.
\[ \log\sqrt{\frac{xz^2}{y^{16}}} \]}
{\( \frac12\log x-8\log y+\log z \)}{\( \frac12\log x+8\log y-\log z \)}{\( 8\log x+\frac12\log y-\log z \)}{\( \log x-16\log y+2\log z \)}{}{}}
\End

\Data\ID{1003102414}
\Updated{2023-04-03 01:04:45}
\Level{B}
\SubArea{Calculations with logarithms}
\LANGen{
\Question{
From the given  expressions choose the one   equivalent  to \( \log\left( 8\cdot\sqrt[3]{75} \right) \), if \( \log2=a\), \( \log3=b \) and \( \log5=c \).}
{\( 3a+\frac13 b+\frac23 c \)}{\( 3a+\frac13 b+\frac13 c \)}{\( 4a+\frac13 b+\frac23 c \)}{\( a+\frac13 b+\frac23 c \)}{}{}}
\LANGpl{
\Question{
Z podanych wyrażeń wybierz to, które jest  odpowiednikiem  \( \log\left( 8\cdot\sqrt[3]{75} \right) \), jeśli \( \log2=a\), \( \log3=b \) i \( \log5=c \).}
{\( 3a+\frac13 b+\frac23 c \)}{\( 3a+\frac13 b+\frac13 c \)}{\( 4a+\frac13 b+\frac23 c \)}{\( a+\frac13 b+\frac23 c \)}{}{}}
\LANGcs{
\Question{
Určete hodnotu výrazu \( \log\left( 8\cdot\sqrt[3]{75} \right) \), když \( \log2=a\), \( \log3=b \) a \( \log5=c \).}
{\( 3a+\frac13 b+\frac23 c \)}{\( 3a+\frac13 b+\frac13 c \)}{\( 4a+\frac13 b+\frac23 c \)}{\( a+\frac13 b+\frac23 c \)}{}{}}
\LANGsk{
\Question{
Určte hodnotu výrazu \( \log\left( 8\cdot\sqrt[3]{75} \right) \), ak \( \log2=a\), \( \log3=b \) a \( \log5=c \).}
{\( 3a+\frac13 b+\frac23 c \)}{\( 3a+\frac13 b+\frac13 c \)}{\( 4a+\frac13 b+\frac23 c \)}{\( a+\frac13 b+\frac23 c \)}{}{}}
\LANGes{
\Question{
Elige una de las expresiones siguientes que sea equivalente a \( \log\left( 8\cdot\sqrt[3]{75} \right) \), si \( \log2=a\), \( \log3=b \) y \( \log5=c \).}
{\( 3a+\frac13 b+\frac23 c \)}{\( 3a+\frac13 b+\frac13 c \)}{\( 4a+\frac13 b+\frac23 c \)}{\( a+\frac13 b+\frac23 c \)}{}{}}
\End

\Data\ID{2000014106}
\Updated{2023-04-03 01:04:37}
\Level{B}
\SubArea{Calculations with logarithms}
\LANGen{
\Question{
Find the value of \( \log_2 5 + \log_2 20\) if \(\log_2 5=a\) and \( \log_2 20=b\).}
{\( 2a+2\)}{\( 2^a+2^b\)}{\( ab\)}{\( 5a\)}{}{}}
\LANGpl{
\Question{
Znaidź wartość \( \log_2 5 + \log_2 20\), jeśli \(\log_2 5=a\) i \( \log_2 20=b\).}
{\( 2a+2\)}{\( 2^a+2^b\)}{\( ab\)}{\( 5a\)}{}{}}
\LANGcs{
\Question{
Určete hodnotu výrazu \( \log_2 5 + \log_2 20\), když \(\log_2 5=a\) a \( \log_2 20=b\).}
{\( 2a+2\)}{\( 2^a+2^b\)}{\( ab\)}{\( 5a\)}{}{}}
\LANGsk{
\Question{
Určte hodnotu výrazu \( \log_2 5 + \log_2 20\), ak \(\log_2 5=a\) a \( \log_2 20=b\).}
{\( 2a+2\)}{\( 2^a+2^b\)}{\( ab\)}{\( 5a\)}{}{}}
\LANGes{
\Question{
Halla el valor de \( \log_2 5 + \log_2 20\) si \(\log_2 5=a\) y \( \log_2 20=b\).}
{\( 2a+2\)}{\( 2^a+2^b\)}{\( ab\)}{\( 5a\)}{}{}}
\End

\Data\ID{2010011005}
\Updated{2023-04-03 01:04:03}
\Level{B}
\SubArea{Calculations with logarithms}
\LANGen{
\Question{
If \( a \), \( b \), \( c\in(0;\infty) \) then the expression \( \log_2a+3 \log_2 b-\frac12 \log_2c \) is equivalent to:}
{\( \log_2\frac{ab^3}{\sqrt{c}} \)}{\( \log_2\frac{3ab}{\frac12 c} \)}{\( \log_2 \left({ab^3}{c}^{\frac12} \right)\)}{\( \log_2 \left(-\frac32 abc\right) \)}{}{}}
\LANGpl{
\Question{
Jeśli \( a \), \( b \), \( c\in(0;\infty) \), to wyrażenie \( \log_2a+3 \log_2 b-\frac12 \log_2c \) jest równoważne:}
{\( \log_2\frac{ab^3}{\sqrt{c}} \)}{\( \log_2\frac{3ab}{\frac12 c} \)}{\( \log_2 \left({ab^3}{c}^{\frac12} \right)\)}{\( \log_2 \left(-\frac32 abc\right) \)}{}{}}
\LANGcs{
\Question{
Když \( a \), \( b \), \( c\in(0;\infty) \), pak výraz \( \log_2a+3 \log_2 b-\frac12 \log_2c \) je ekvivalentní s výrazem:}
{\( \log_2\frac{ab^3}{\sqrt{c}} \)}{\( \log_2\frac{3ab}{\frac12 c} \)}{\( \log_2 \left({ab^3}{c}^{\frac12} \right)\)}{\( \log_2 \left(-\frac32 abc\right) \)}{}{}}
\LANGsk{
\Question{
Nech \( a \), \( b \), \( c\in(0;\infty) \). Daný výraz \[ \log_2a+3 \log_2 b-\frac12 \log_2c \] sa rovná:}
{\( \log_2\frac{ab^3}{\sqrt{c}} \)}{\( \log_2\frac{3ab}{\frac12 c} \)}{\( \log_2 \left({ab^3}{c}^{\frac12} \right)\)}{\( \log_2 \left(-\frac32 abc\right) \)}{}{}}
\LANGes{
\Question{
Si \( a \), \( b \), \( c\in(0;\infty) \) la expresión \( \log_2a+3 \log_2 b-\frac12 \log_2c \) es equivalente a:}
{\( \log_2\frac{ab^3}{\sqrt{c}} \)}{\( \log_2\frac{3ab}{\frac12 c} \)}{\( \log_2 \left({ab^3}{c}^{\frac12} \right)\)}{\( \log_2 \left(-\frac32 abc\right) \)}{}{}}
\End

\Data\ID{2010016001}
\Updated{2023-04-03 01:04:37}
\Level{B}
\SubArea{Calculations with logarithms}
\LANGen{
\Question{
The value of the expression \( \log 25^4+\log 4^4\) is:}
{\(8\)}{\( 4\)}{\( 4+\log 29\)}{\( 8+\log 29\)}{}{}}
\LANGpl{
\Question{
Wartość wyrażenia \( \log 25^4+\log 4^4\) jest równa:}
{\(8\)}{\( 4\)}{\( 4+\log 29\)}{\( 8+\log 29\)}{}{}}
\LANGcs{
\Question{
Hodnota výrazu \( \log 25^4+\log 4^4\) je:}
{\(8\)}{\( 4\)}{\( 4+\log 29\)}{\( 8+\log 29\)}{}{}}
\LANGsk{
\Question{
Hodnota výrazu \( \log 25^4+\log 4^4\) sa rovná číslu:}
{\(8\)}{\( 4\)}{\( 4+\log 29\)}{\( 8+\log 29\)}{}{}}
\LANGes{
\Question{
El valor de la expresión \( \log 25^4+\log 4^4\) es:}
{\(8\)}{\( 4\)}{\( 4+\log 29\)}{\( 8+\log 29\)}{}{}}
\End

\Data\ID{2010016002}
\Updated{2023-04-03 01:04:56}
\Level{B}
\SubArea{Calculations with logarithms}
\LANGen{
\Question{
The value of the expression \( \log 125^2+\log 8^2\) is:}
{\(6\)}{\( 5\)}{\( 4+\log 133\)}{\( 8+\log 133\)}{}{}}
\LANGpl{
\Question{
Wartość wyrażenia \( \log 125^2+\log 8^2\) jest równa:}
{\(6\)}{\( 5\)}{\( 4+\log 133\)}{\( 8+\log 133\)}{}{}}
\LANGcs{
\Question{
Hodnota výrazu \( \log 125^2+\log 8^2\) je:}
{\(6\)}{\( 5\)}{\( 4+\log 133\)}{\( 8+\log 133\)}{}{}}
\LANGsk{
\Question{
Hodnota výrazu \( \log 125^2+\log 8^2\) sa rovná číslu:}
{\(6\)}{\( 5\)}{\( 4+\log 133\)}{\( 8+\log 133\)}{}{}}
\LANGes{
\Question{
El valor de la expresión \( \log 125^2+\log 8^2\) es:}
{\(6\)}{\( 5\)}{\( 4+\log 133\)}{\( 8+\log 133\)}{}{}}
\End

\Data\ID{2010016003}
\Updated{2023-04-03 01:04:59}
\Level{B}
\SubArea{Calculations with logarithms}
\LANGen{
\Question{
The number \( 1+\log_2 7\) equals:}
{\(\log_2 14\)}{\( 3\)}{\( 4.5\)}{\( \log_2 9\)}{}{}}
\LANGpl{
\Question{
Liczba \( 1+\log_2 7\) jest równa:}
{\(\log_2 14\)}{\( 3\)}{\( 4{,}5\)}{\( \log_2 9\)}{}{}}
\LANGcs{
\Question{
Číslo \( 1+\log_2 7\) je rovno:}
{\(\log_2 14\)}{\( 3\)}{\( 4{,}5\)}{\( \log_2 9\)}{}{}}
\LANGsk{
\Question{
Číslo \( 1+\log_2 7\) sa rovná:}
{\(\log_2 14\)}{\( 3\)}{\( 4{,}5\)}{\( \log_2 9\)}{}{}}
\LANGes{
\Question{
El número \( 1+\log_2 7\) es igual a:}
{\(\log_2 14\)}{\( 3\)}{\( 4.5\)}{\( \log_2 9\)}{}{}}
\End

\Data\ID{2010016004}
\Updated{2023-04-03 01:04:19}
\Level{B}
\SubArea{Calculations with logarithms}
\LANGen{
\Question{
The number \( \log_3 6 -1\) equals:}
{\(\log_3 2\)}{\( \log_3 5\)}{\( 5\)}{\( 1\)}{}{}}
\LANGpl{
\Question{
Liczba \( \log_3 6 -1\) jest równa:}
{\(\log_3 2\)}{\( \log_3 5\)}{\( 5\)}{\( 1\)}{}{}}
\LANGcs{
\Question{
Číslo \( \log_3 6 -1\) je rovno:}
{\(\log_3 2\)}{\( \log_3 5\)}{\( 5\)}{\( 1\)}{}{}}
\LANGsk{
\Question{
Číslo \( \log_3 6 -1\) sa rovná:}
{\(\log_3 2\)}{\( \log_3 5\)}{\( 5\)}{\( 1\)}{}{}}
\LANGes{
\Question{
El número \( \log_3 6 -1\) es igual a:}
{\(\log_3 2\)}{\( \log_3 5\)}{\( 5\)}{\( 1\)}{}{}}
\End

\Data\ID{1003102407}
\Updated{2023-04-03 01:04:30}
\Level{C}
\SubArea{Calculations with logarithms}
\LANGen{
\Question{
If \( \log_ab=100 \) and \( \log_4a=10 \) then the value of \( b \) is:}
{\( 2^{2000} \)}{\( 2^{1000} \)}{\( 2^{1002} \)}{\( 2^{200} \)}{}{}}
\LANGpl{
\Question{
Jeśli \( \log_ab=100 \) i \( \log_4a=10 \) wtedy wartość \( b \) wynosi:}
{\( 2^{2000} \)}{\( 2^{1000} \)}{\( 2^{1002} \)}{\( 2^{200} \)}{}{}}
\LANGcs{
\Question{
Když platí \( \log_ab=100 \) a \( \log_4a=10 \), tak hodnota \( b \) se rovná:}
{\( 2^{2000} \)}{\( 2^{1000} \)}{\( 2^{1002} \)}{\( 2^{200} \)}{}{}}
\LANGsk{
\Question{
Ak platí \( \log_ab=100 \) a \( \log_4a=10 \), tak hodnota \( b \) sa rovná:}
{\( 2^{2000} \)}{\( 2^{1000} \)}{\( 2^{1002} \)}{\( 2^{200} \)}{}{}}
\LANGes{
\Question{
Si \( \log_ab=100 \) y \( \log_4a=10 \), el valor de \( b \) es:}
{\( 2^{2000} \)}{\( 2^{1000} \)}{\( 2^{1002} \)}{\( 2^{200} \)}{}{}}
\End

\Data\ID{1003102410}
\Updated{2023-04-03 01:04:48}
\Level{C}
\SubArea{Calculations with logarithms}
\LANGen{
\Question{
If \( x\in(0;1)\cup(1;\infty) \) then the product \( \left(\log_x3\right)\left(\log_5x\right) \) can be written as:}
{\( \log_53 \)}{\( \log_35 \)}{\( \log_{5x}(3x) \)}{\( \log_x3+\log_5x \)}{}{}}
\LANGpl{
\Question{
Jeśli \( x\in(0;1)\cup(1;\infty) \), wtedy iloczyn \( \left(\log_x3\right)\left(\log_5x\right) \) może zostać zapisany jako:}
{\( \log_53 \)}{\( \log_35 \)}{\( \log_{5x}(3x) \)}{\( \log_x3+\log_5x \)}{}{}}
\LANGcs{
\Question{
Když \( x\in(0;1)\cup(1;\infty) \), tak součin \( \left(\log_x3\right)\left(\log_5x\right) \) můžeme zapsat ve tvaru:}
{\( \log_53 \)}{\( \log_35 \)}{\( \log_{5x}(3x) \)}{\( \log_x3+\log_5x \)}{}{}}
\LANGsk{
\Question{
Ak \( x\in(0;1)\cup(1;\infty) \), tak súčin \( \left(\log_x3\right)\left(\log_5x\right) \) môžme zapísať v tvare:}
{\( \log_53 \)}{\( \log_35 \)}{\( \log_{5x}(3x) \)}{\( \log_x3+\log_5x \)}{}{}}
\LANGes{
\Question{
Si \( x\in(0;1)\cup(1;\infty) \), el producto \( \left(\log_x3\right)  \left(\log_5x\right) \) se puede escribir como:}
{\( \log_53 \)}{\( \log_35 \)}{\( \log_{5x}(3x) \)}{\( \log_x3+\log_5x \)}{}{}}
\End

\Data\ID{1003102415}
\Updated{2023-04-03 01:04:03}
\Level{C}
\SubArea{Calculations with logarithms}
\LANGen{
\Question{
Let \( a \in(0;\infty) \). The expression \( \log_4a-\log_{16}a-\log_{\frac14}a \) is equivalent to:}
{\( \frac32 \log_4a \)}{\( -\frac12\log_4a \)}{\( \log_4a \)}{\( 0 \)}{}{}}
\LANGpl{
\Question{
Niech  \( a \in(0;\infty) \). Wyrażenie \( \log_4a-\log_{16}a-\log_{\frac14}a \) jest równe:}
{\( \frac32 \log_4a \)}{\( -\frac12\log_4a \)}{\( \log_4a \)}{\( 0 \)}{}{}}
\LANGcs{
\Question{
Nechť \( a \in(0;\infty) \). Výraz \( \log_4a-\log_{16}a-\log_{\frac14}a \) se rovná:}
{\( \frac32 \log_4a \)}{\( -\frac12\log_4a \)}{\( \log_4a \)}{\( 0 \)}{}{}}
\LANGsk{
\Question{
Nech \( a \in(0;\infty) \). Výraz \( \log_4a-\log_{16}a-\log_{\frac14}a \) sa rovná:}
{\( \frac32 \log_4a \)}{\( -\frac12\log_4a \)}{\( \log_4a \)}{\( 0 \)}{}{}}
\LANGes{
\Question{
Si \( a \in(0;\infty) \), la expresión \( \log_4a-\log_{16}a-\log_{\frac14}a \) es equivalente a:}
{\( \frac32 \log_4a \)}{\( -\frac12\log_4a \)}{\( \log_4a \)}{\( 0 \)}{}{}}
\End

\Data\ID{2000014103}
\Updated{2023-04-03 01:04:33}
\Level{C}
\SubArea{Calculations with logarithms}
\LANGen{
\Question{
Find the value of \( \log_{ab}x\) if \(\log_a x=2\) and \(\log_b x=3\).}
{\( \frac65\)}{\( \frac16\)}{\( 6\)}{\( \frac56\)}{}{}}
\LANGpl{
\Question{
Znajdź wartość \( \log_{ab}x\) jeśli \(\log_a x=2\) i \(\log_b x=3\).}
{\( \frac65\)}{\( \frac16\)}{\( 6\)}{\( \frac56\)}{}{}}
\LANGcs{
\Question{
Určete hodnotu výrazu \( \log_{ab}x\), kdy \(\log_a x=2\) a \(\log_b x=3\).}
{\( \frac65\)}{\( \frac16\)}{\( 6\)}{\( \frac56\)}{}{}}
\LANGsk{
\Question{
Určte hodnotu výrazu \( \log_{ab}x\), ak \(\log_a x=2\) a \(\log_b x=3\).}
{\( \frac65\)}{\( \frac16\)}{\( 6\)}{\( \frac56\)}{}{}}
\LANGes{
\Question{
Halla el valor de \( \log_{ab}x\) si \(\log_a x=2\) y \(\log_b x=3\).}
{\( \frac65\)}{\( \frac16\)}{\( 6\)}{\( \frac56\)}{}{}}
\End

\Data\ID{2000014104}
\Updated{2023-04-03 01:04:49}
\Level{C}
\SubArea{Calculations with logarithms}
\LANGen{
\Question{
Find the value of \(L\) if \(L=\log_{\sqrt{2}}2 \cdot \log_2 \sqrt{3} \cdot \log_{\sqrt{3}} 4\).}
{\( L=4\)}{\( L=2\)}{\( L=3\)}{\( L=1\)}{}{}}
\LANGpl{
\Question{
Znajdź wartość \(L\), jeśli \(L=\log_{\sqrt{2}}2 \cdot \log_2 \sqrt{3} \cdot \log_{\sqrt{3}} 4\).}
{\( L=4\)}{\( L=2\)}{\( L=3\)}{\( L=1\)}{}{}}
\LANGcs{
\Question{
Určete hodnotu \(L\), když \(L=\log_{\sqrt{2}}2 \cdot \log_2 \sqrt{3} \cdot \log_{\sqrt{3}} 4\).}
{\( L=4\)}{\( L=2\)}{\( L=3\)}{\( L=1\)}{}{}}
\LANGsk{
\Question{
Určte hodnotu \(L\), ak \(L=\log_{\sqrt{2}}2 \cdot \log_2 \sqrt{3} \cdot \log_{\sqrt{3}} 4\).}
{\( L=4\)}{\( L=2\)}{\( L=3\)}{\( L=1\)}{}{}}
\LANGes{
\Question{
Halla el valor de \(L\) si \(L=\log_{\sqrt{2}}2 \cdot \log_2 \sqrt{3} \cdot \log_{\sqrt{3}} 4\).}
{\( L=4\)}{\( L=2\)}{\( L=3\)}{\( L=1\)}{}{}}
\End

\Data\ID{2000014105}
\Updated{2023-04-03 01:04:33}
\Level{C}
\SubArea{Calculations with logarithms}
\LANGen{
\Question{
Find the value of \( \left( \log_{\frac1a}b\right) \cdot \left( \log_{\frac1b}c\right) \cdot \left( \log_{\frac1c}a\right)\).}
{\( -1\)}{\( \frac1{abc}\)}{\( abc\)}{\( 1\)}{}{}}
\LANGpl{
\Question{
Znajdź wartość \( \left( \log_{\frac1a}b\right) \cdot \left( \log_{\frac1b}c\right) \cdot \left( \log_{\frac1c}a\right)\).}
{\( -1\)}{\( \frac1{abc}\)}{\( abc\)}{\( 1\)}{}{}}
\LANGcs{
\Question{
Určete hodnotu výrazu \( \left( \log_{\frac1a}b\right) \cdot \left( \log_{\frac1b}c\right) \cdot \left( \log_{\frac1c}a\right)\).}
{\( -1\)}{\( \frac1{abc}\)}{\( abc\)}{\( 1\)}{}{}}
\LANGsk{
\Question{
Určte hodnotu výrazu \( \left( \log_{\frac1a}b\right) \cdot \left( \log_{\frac1b}c\right) \cdot \left( \log_{\frac1c}a\right)\).}
{\( -1\)}{\( \frac1{abc}\)}{\( abc\)}{\( 1\)}{}{}}
\LANGes{
\Question{
Halla el valor de \( \left( \log_{\frac1a}b\right) \cdot \left( \log_{\frac1b}c\right) \cdot \left( \log_{\frac1c}a\right)\).}
{\( -1\)}{\( \frac1{abc}\)}{\( abc\)}{\( 1\)}{}{}}
\End

\Data\ID{2000014107}
\Updated{2023-04-03 01:04:49}
\Level{C}
\SubArea{Calculations with logarithms}
\LANGen{
\Question{
Choose the correct equality.}
{\( 4^{\log_23}=9\)}{\( 2^{1-\log_23}=3\)}{\( 4^{\log_24}=4\)}{\( 4^{1+\log_42}=16\)}{}{}}
\LANGpl{
\Question{
Wybierz prawdziwą równość.}
{\( 4^{\log_23}=9\)}{\( 2^{1-\log_23}=3\)}{\( 4^{\log_24}=4\)}{\( 4^{1+\log_42}=16\)}{}{}}
\LANGcs{
\Question{
Vyberte pravdivé tvrzení.}
{\( 4^{\log_23}=9\)}{\( 2^{1-\log_23}=3\)}{\( 4^{\log_24}=4\)}{\( 4^{1+\log_42}=16\)}{}{}}
\LANGsk{
\Question{
Vyberte pravdivé tvrdenie.}
{\( 4^{\log_23}=9\)}{\( 2^{1-\log_23}=3\)}{\( 4^{\log_24}=4\)}{\( 4^{1+\log_42}=16\)}{}{}}
\LANGes{
\Question{
Elige la igualdad correcta.}
{\( 4^{\log_23}=9\)}{\( 2^{1-\log_23}=3\)}{\( 4^{\log_24}=4\)}{\( 4^{1+\log_42}=16\)}{}{}}
\End

\Data\ID{2000014108}
\Updated{2023-04-03 01:04:10}
\Level{C}
\SubArea{Calculations with logarithms}
\LANGen{
\Question{
Find the value of \(\log_{y^2x}y^7x^5\) if \(\log_x y=-2\).}
{\( 3\)}{\( -3\)}{\( 17.5\)}{\( \frac{17}3\)}{}{}}
\LANGpl{
\Question{
Znajdż wartość \(\log_{y^2x}y^7x^5\), jeśli \(\log_x y=-2\).}
{\( 3\)}{\( -3\)}{\( 17{,}5\)}{\( \frac{17}3\)}{}{}}
\LANGcs{
\Question{
Určete hodnotu výrazu \(\log_{y^2x}y^7x^5\), když \(\log_x y=-2\).}
{\( 3\)}{\( -3\)}{\( 17{,}5\)}{\( \frac{17}3\)}{}{}}
\LANGsk{
\Question{
Určte hodnotu výrazu \(\log_{y^2x}y^7x^5\), ak \(\log_x y=-2\).}
{\( 3\)}{\( -3\)}{\( 17{,}5\)}{\( \frac{17}3\)}{}{}}
\LANGes{
\Question{
Halla el valor de \(\log_{y^2x}y^7x^5\) si \(\log_x y=-2\).}
{\( 3\)}{\( -3\)}{\( 17.5\)}{\( \frac{17}3\)}{}{}}
\End

\Data\ID{2010011004}
\Updated{2023-04-03 01:04:11}
\Level{C}
\SubArea{Calculations with logarithms}
\LANGen{
\Question{
If \( x\in(0;1)\cup(1;\infty) \) then the product \( \left(\log_x4\right)\left(\log_{16}x\right) \) can be written as:}
{\( \frac12 \)}{\( 2 \)}{\( \log_x 4 + \log_{16} x \)}{\( \frac1{4} \)}{}{}}
\LANGpl{
\Question{
Jeżeli \( x\in(0;1)\cup(1;\infty) \), to iloczyn \( \left(\log_x4\right)\left(\log_{16}x\right) \) można zapisać jako:}
{\( \frac12 \)}{\( 2 \)}{\( \log_x 4 + \log_{16} x \)}{\( \frac1{4} \)}{}{}}
\LANGcs{
\Question{
Když \( x\in(0;1)\cup(1;\infty) \), pak součin \( \left(\log_x4\right)\left(\log_{16}x\right) \) lze zapsat ve tvaru:}
{\( \frac12 \)}{\( 2 \)}{\( \log_x 4 + \log_{16} x \)}{\( \frac1{4} \)}{}{}}
\LANGsk{
\Question{
Ak \( x\in(0;1)\cup(1;\infty) \), tak súčin \( \left(\log_x4\right)\left(\log_{16}x\right) \) môžme zapísať v tvare:}
{\( \frac12 \)}{\( 2 \)}{\( \log_x 4 + \log_{16} x \)}{\( \frac1{4} \)}{}{}}
\LANGes{
\Question{
Si \( x\in(0;1)\cup(1;\infty) \) el producto \( \left(\log_x4\right)\left(\log_{16}x\right) \) se puede escribir como:}
{\( \frac12 \)}{\( 2 \)}{\( \log_x 4 + \log_{16} x \)}{\( \frac1{4} \)}{}{}}
\End
